\documentclass[../Main.tex]{subfiles}
\begin{document}

\label{sec:methodology}

This section outlines the technology stack and development tools employed in the implementation of the system. The selection of these technologies is based on criteria such as performance, scalability, community support, and suitability for the specific requirements of a flood relief and charity management application.

\section{Frontend Development: Flutter}
For the mobile application development, Flutter, an open-source UI software development kit created by Google, was selected.

\textbf{Key Strengths:}
\begin{itemize}
    \item[] \textbf{Cross-Platform Capabilities:} Flutter allows the creation of natively compiled applications for mobile, web, and desktop from a single codebase (Dart language). This significantly reduces development time compared to maintaining separate codebases for Android and iOS.
    \item[] \textbf{High Performance:} Unlike other frameworks that rely on a Javascript bridge, Flutter uses the Skia graphics engine to render UI components directly, resulting in smooth animations (60fps) and native-like performance.
    \item[] \textbf{Hot Reload:} This feature enables developers to see changes in the code instantly without restarting the app, accelerating the debugging process.
\end{itemize}

\textbf{Comparison with React Native:}
While React Native is a popular alternative, Flutter provides a more consistent UI across platforms because it does not rely on OEM widgets. Furthermore, Flutter's performance is generally superior in complex animation scenarios due to the absence of the JavaScript bridge bottleneck.

\section{Backend Development: NestJS}
The server-side logic is built using NestJS, a progressive Node.js framework for building efficient, reliable, and scalable server-side applications.

\textbf{Key Strengths:}
\begin{itemize}
    \item[] \textbf{Modular Architecture:} NestJS is built around the concept of modules, controllers, and services (inspired by Angular). This structure forces code organization, making the system easier to maintain and scale.
    \item[] \textbf{TypeScript Support:} It fully supports TypeScript, providing strong typing which helps prevent runtime errors and improves code readability.
    \item[] \textbf{Dependency Injection:} Its built-in Dependency Injection (DI) container promotes loose coupling and enhances testability.
\end{itemize}

\textbf{Comparison with Express.js:}
Express.js is a minimalist and unopinionated framework, which often leads to "spaghetti code" in large projects as developers must define their own architecture. NestJS, conversely, provides an "opinionated" structure out-of-the-box, which is crucial for maintaining the complexity of a system handling financial transactions (donations) and real-time rescue operations.

\section{Database Management System: PostgreSQL}
PostgreSQL is utilized as the primary relational database management system (RDBMS) for storing user data, campaign details, and transaction logs.

\textbf{Key Strengths:}
\begin{itemize}
    \item[] \textbf{ACID Compliance:} PostgreSQL ensures validity even in the event of errors or power failures, which is critical for financial transactions in the donation module.
    \item[] \textbf{Advanced Data Types:} It supports JSON/JSONB natively, allowing for a hybrid approach (relational + NoSQL) when handling flexible data structures.
    \item[] \textbf{Geospatial Support (PostGIS):} Although not explicitly detailed here, PostgreSQL's extensibility supports complex geospatial queries, which is beneficial for location-based services in rescue operations.
\end{itemize}

\textbf{Comparison with MySQL:}
While MySQL is widely used for web applications, PostgreSQL is generally considered more robust for complex queries and write-heavy operations. Its strict adherence to SQL standards and superior handling of concurrent transactions make it a better choice for this system's reliability requirements.

\section{Object-Relational Mapping: Prisma}
To interact with the PostgreSQL database, the system utilizes **Prisma**, a next-generation open-source ORM. It acts as the data access layer, connecting the NestJS backend to the database.

\textbf{Key Strengths:}
\begin{itemize}
    \item[] \textbf{Type Safety:} Prisma automatically generates a strongly-typed TypeScript client based on the schema definition. This ensures that database queries are validated at compile-time, significantly reducing runtime errors caused by data type mismatches.
    \item[] \textbf{Schema-First Approach:} The database structure is defined in a centralized, human-readable file (`schema.prisma`). This file serves as the single source of truth, making it easy to model data and manage database migrations.
    \item[] \textbf{Intuitive API:} Prisma's query syntax is designed to be readable and concise, simplifying complex operations like nested writes or filtering deep relationships.
\end{itemize}

\textbf{Comparison with TypeORM:}
TypeORM is the traditional default for NestJS and relies heavily on classes and decorators. While powerful, TypeORM can be verbose and prone to synchronization issues between the entity classes and the actual database state. Prisma outperforms TypeORM in terms of Developer Experience (DX) by offering superior auto-completion (IntelliSense) and a more predictable query engine that efficiently handles complex data relations, often avoiding performance pitfalls like the "N+1 problem" more effectively.

\section{Communication Protocols}
To ensure seamless communication between the client, server, and IoT devices (if any), the system employs a hybrid protocol approach.

\subsection{HTTP/HTTPS}
Standard HTTP/HTTPS protocols are used for RESTful API communication (e.g., login, fetching campaign lists, updating user profiles). HTTPS ensures that sensitive data, such as authentication tokens and personal information, is encrypted during transmission.

\subsection{MQTT (Message Queuing Telemetry Transport)}
Given the requirement to operate in potentially weak network environments during disasters, MQTT is implemented for critical updates and location sharing.

\textbf{Key Strengths:}
\begin{itemize}
    \item[] \textbf{Lightweight:} MQTT header size is minimal compared to HTTP headers, reducing bandwidth consumption.
    \item[] \textbf{Reliability:} It supports Quality of Service (QoS) levels, ensuring message delivery even in unstable network conditions.
    \item[] \textbf{Pub/Sub Model:} This decouples the sender and receiver, allowing for efficient real-time data distribution (e.g., broadcasting emergency alerts).
\end{itemize}

\section{MQTT Broker: Eclipse Mosquitto}
To facilitate real-time data transmission during the development and testing phases, the system utilizes Eclipse Mosquitto, an open-source message broker that implements the MQTT protocol versions 5.0, 3.1.1, and 3.1.

\textbf{Key Strengths:}
\begin{itemize}
    \item[] \textbf{Lightweight and Efficient:} Written in C, Mosquitto is extremely lightweight and consumes minimal resources (CPU and RAM). This makes it ideal for running on local development machines or low-power IoT gateways (edge devices) without causing system lag.
    \item[] \textbf{Standard Compliance:} It strictly adheres to MQTT standards, ensuring that the client implementations (Flutter app and backend services) remain compatible with other brokers if the system needs to switch in the future.
    \item[] \textbf{Ease of Deployment:} Mosquitto provides official Docker images, allowing developers to spin up a fully functional broker with a single command (`docker run`), drastically reducing the time required for environment setup.
\end{itemize}

\textbf{Role in Development Environment:}
In this project, Mosquitto serves as the central hub for testing the "Publish/Subscribe" mechanism. It allows the development team to simulate unstable network conditions and verify that the IoT devices and the mobile application can successfully reconnect and sync data (e.g., location coordinates, emergency alerts) without data loss.

\textbf{Comparison with EMQX / HiveMQ:}
While **Mosquitto** is excellent for development and edge computing due to its simplicity, **EMQX** or **HiveMQ** are often preferred for large-scale production environments requiring clustering and massive concurrency (millions of connections). However, for the scope of this project and the current development stage, Mosquitto offers the best balance between performance and ease of use, avoiding the configuration overhead of complex enterprise brokers.

\section{Containerization: Docker}
**Docker** is employed to containerize the application, packaging the code along with its dependencies into standardized units for development and deployment.

\textbf{Key Strengths:}
\begin{itemize}
    \item[] \textbf{Environment Consistency:} Docker eliminates the common "it works on my machine" issue. By bundling the application with its specific runtime environment (Node.js version, OS libraries), it guarantees that the software runs exactly the same way on a developer's laptop as it does on the production server.
    \item[] \textbf{Isolation:} The backend, database, and other services run in isolated containers. This prevents conflicts between different projects or dependency versions on the same host machine.
    \item[] \textbf{Rapid Deployment:} Using `docker-compose`, the entire system stack (NestJS API, PostgreSQL database, ...) can be orchestrated and launched with a single command, streamlining the setup process for new developers.
\end{itemize}

\textbf{Comparison with Virtual Machines (VMs):}
Unlike traditional Virtual Machines that require a full Operating System for each instance, Docker containers share the host system's kernel. This makes containers significantly more lightweight (occupying MBs instead of GBs) and much faster to start up (seconds vs. minutes), resulting in better resource efficiency and scalability for the application.

\section{Third-Party Services and APIs}
\begin{itemize}
    \item[] \textbf{Google Authentication:} Integrated to provide a secure and convenient login mechanism (SSO), reducing the friction of user registration.
    \item[] \textbf{Google Maps API:} Used to visualize the rescue map, precipitation heatmaps, and route navigation. Google Maps is the industry standard for accuracy and detail.
    \item[] \textbf{Webhooks:} The system utilizes webhooks to receive real-time updates from external banking systems. This eliminates the need for constant polling (checking the server repeatedly), thereby saving server resources and ensuring immediate transaction verification.
\end{itemize}

\section{Development and Testing Tools}
\begin{itemize}
    \item[] \textbf{IDE: Visual Studio Code (VS Code):} Selected for its lightweight nature and vast ecosystem of extensions (for Dart, TypeScript, Docker, etc.). Its IntelliSense and debugging tools significantly enhance developer productivity.
    \item[] \textbf{API Testing: Postman:} Used to design, document, and test API endpoints. Postman allows developers to simulate requests, inspect responses, and automate testing scenarios before integrating APIs into the frontend.
\end{itemize}

\end{document}